\clearpage
\fancyhead[L]{\small CARLINI. MILAN EPHEMERIDES FOR 1816.}
\SourceLeaf{597}
\newcommand{\ObservedLibelleDenominator}{\hbox{\vrule width0pt height4.600000bp depth0pt\special{pdf:bcontent}\special{pdf:code q 0 g 2.146667 4.600000 m 2.146667 4.293333 l 2.453333 3.986667 l 2.453333 3.680000 l 1.840000 3.066667 l 1.533333 3.066667 l 0.613333 2.146667 l 0.000000 2.146667 l 0.000000 0.613333 l 0.306667 0.613333 l 0.613333 0.306667 l 1.226667 0.306667 l 1.533333 0.000000 l 3.066667 0.000000 l 3.986667 0.920000 l 3.986667 1.840000 l 3.066667 2.760000 l 3.066667 3.066667 l 3.680000 3.680000 l 3.680000 3.986667 l 3.066667 4.600000 l h 1.533333 2.146667 m 1.840000 2.453333 l 2.453333 1.840000 l 2.453333 1.533333 l 1.840000 0.920000 l 1.533333 1.226667 l h f* Q}\special{pdf:econtent}\kern3.986667bp}}

\begin{center}\small Göttingische gelehrte Anzeigen. No. 47. Pages 457..46c. 23 March 1816.\\[3pt]\rule{.72\linewidth}{.3pt}\end{center}

\emph{Effemeridi astronomiche di Milano per l’anno bisestile} 1816, \emph{calcolate da} \textsc{Francesco Carlini}. \emph{Con appendice. Dalla cesarea regia stamperia.} Milan, 1815. Octavo. The ephemerides comprise 108 pages, the appendix 100.

The supplements accompanying the astronomical calendar this time are 1) the list of the 34 fundamental stars, together with their annual precessions and the constants for calculating their aberration and nutation, as in the 1815 volume (see our Anzeigen for last year, p. 89 [p. 574 of this volume]). 2) \textsc{Carlini}’s refraction table. This table, based on \textsc{Laplace}’s hypothesis about the state of the atmosphere and on the theory of refraction derived from it, but with the constants involved newly determined from his own observations, had already been printed in the 1808 volume of the ephemerides. As the reviewer does not have that volume to hand, he cannot say whether the present table is simply an unchanged reprint. 3) A list of occultations of fixed stars in 1816, calculated by the Florentine astronomers for the meridian and latitude of Florence.

The appendix contains the following papers. 1) Observed refraction at low altitudes above the horizon, by \textsc{Barnabas Oriani}. The observations with \textsc{Reichenbach}’s three-foot repeating circle published in earlier volumes include 19 lower culminations of Capella observed in 1811. In Milan this star transits at an altitude of $1^\circ\;36'$, only a few minutes of arc above the mountains that bound the horizon. In addition, in 1811 a further 14 upper culminations had been observed, from which, using the latitude $45^\circ\;28'\;0''7$, the mean declination $45^\circ\;47'\;28''5$ for the beginning of 1811 was derived. Mr.\ \textsc{Oriani} calculated the lower culminations using the principal refraction tables, namely those of \textsc{Bradley}, \textsc{Tob. Mayer}, \textsc{Piazzi}, \textsc{Delambre}, \textsc{Carlini} and \textsc{Bessel}. Of these, \textsc{Carlini}’s gave the best agreement, to within a few seconds of arc, provided, that is, that the supposed correction for the northern half of the meridian, which \textsc{Carlini} had earlier thought necessary, was omitted. 2) Continuation of the observations on the periodic swaying of buildings, by \textsc{Angelo Cesaris}. The findings on this subject already reported in the 1813 volume of the ephemerides (see our Anzeigen for 1813, p. 951 [p. 556 of this volume]) are supplemented here by others that are no less remarkable. That account dealt with a movement that could be regarded as something like the rotation, about a vertical axis, of the walls from which the mural quadrant and the transit telescope are suspended. The new findings concern a rotation about a horizontal axis running from east to west. This movement was made visible with an excellent \textsc{Reichenbach} spirit level mounted on the mural quadrant, so sensitive that a change of one second of arc produced a deflection of $1\frac{2}{\ObservedLibelleDenominator}$ lines. The same changes that appeared in this spirit level could also be detected in the zenith distance observed with the mural quadrant and in a second spirit level mounted on the opposite side of the wall. The findings reported here point very clearly to a connection with the weather and with how strongly sunlight warms the building. The changes from morning to afternoon are the most regular. In clear weather the bubble in the spirit level moves southward. The movement averages about $2''$ and disappears entirely when the sky is overcast. 3) Table for the equation of the centre of Vesta (eccentricity $=0.0889$) and \SourceLeaf{598}for the reduction to the ecliptic (inclination of the orbit $=7^\circ\;8'\;20''$, by \textsc{Carlini}. 4) Occultations of stars, observed in Milan in the years 1811--1815 by him, and in Florence in the years 1810--1813 by the astronomers there. --- 5) Obliquity of the ecliptic from solstice observations made with \textsc{Reichenbach}’s three-foot circle, by \textsc{Barnabas Oriani}. The observations at the winter solstice of 1810 and at the two solstices of 1811 had already been reported in earlier volumes. Here we have the detailed observations from the next five solstices and the results drawn from them all. It is very interesting now to have the verdict of \textsc{Reichenbach}’s larger repeating circles as well on this much-discussed subject. The four winter solstices give the mean obliquity of the ecliptic, reduced to the beginning of 1812, $=23^\circ\;27'\;48''20$, and the four summer solstices give $23^\circ\;27'\;50''77$. There is thus still a difference, although much smaller than that found with \textsc{Reichenbach}’s 12-inch circles and other instruments. \textsc{Carlini}’s refraction tables were used to reduce the observations. With \textsc{Delambre}’s refraction tables the difference would be somewhat larger and would amount to $3''71$. --- 6) Meteorological observations in 1814, by \textsc{Angelo Cesaris}, together with a survey of the rainfall in the years 1764--1814. The mean amount of rain over all 51 years comes to 35 inches 3.92 lines. The amounts in individual years naturally vary very irregularly. If they are taken together over a longer series of years, however, there is a steady increase overall. \textsc{Cesaris} attributes this to the increasingly widespread irrigation of fields. In his view, spreading a large quantity of water over a much larger surface exposes it to greater evaporation, so that the water must return as more frequent rain. The number of years is probably still too small, however, and the outcome too dependent on chance fluctuations to draw firm general conclusions. The last year, 1814, was exceptionally wet, with 58 inches 11.58 lines of rain. The smallest amount, in 1771, was 25 inches 11.6 lines.
